\begin{document}

\section{結果}

\subsection{全体結果}

表\ref{tab:overall}に，全体結果を示す．構造化出力を満たした範囲では，\code{gpt-5.4-mini}が\code{gpt-4.1-mini}を大きく上回った．action recallは0.466から0.798へ，critical evidence recallは0.195から0.703へ，orderは0.201から0.553へ上昇し，overclaimは1190から651へ減少した．これは，モデル更新により単なる説明生成だけでなく，証跡接続と順序構成が改善したことを示す．

\begin{table*}[t]
\centering
\scriptsize
\caption{23-chain最終比較の全体結果}
\label{tab:overall}
\begin{tabular}{lrrrrrr}
\toprule
Model & Runs & Action & Crit. ev. & Order & Claim prec. & Overclaim \\
\midrule
\code{gpt-4.1-mini} & 207 & 0.466 & 0.195 & 0.201 & 0.366 & 1190 \\
\code{gpt-5.4-mini} & 207 & 0.798 & 0.703 & 0.553 & 0.584 & 651 \\
\code{gpt-5.5 low raw} & 207 & 0.918 & 0.906 & 0.868 & 0.640 & 1235 \\
\bottomrule
\end{tabular}
\end{table*}

\code{gpt-5.5 low raw}はaction recall 0.918，critical evidence recall 0.906，order 0.868と最も高い復元能力を示した．一方で，overclaimは1235と最大であり，かつ出力契約に失敗した．したがって，\code{gpt-5.5 low raw}は「潜在的な行動復元能力の上限」を示す結果として扱うべきであり，そのまま最良の運用設定と結論づけるべきではない．SOC支援では，高いrecallだけでなく，指定スキーマに従うこと，主系列と近傍証跡を分けること，過剰主張を抑えることが必要である．

\subsection{Stage別結果}

表\ref{tab:bystage}にStage別結果を示す．\code{gpt-5.4-mini}はStage 1からStage 3までaction recallが0.791，0.800，0.802とほぼ一定であり，critical evidence recallはStage 3で0.779と最も高い．これは，alert summary rowsを隠しても，残存telemetryから多くの行動と証跡を復元できることを示す．\code{gpt-5.5 low raw}も全Stageでaction/evidence recallが0.90以上であり，Stage差は小さい．

\begin{table*}[t]
\centering
\scriptsize
\caption{Stage別結果}
\label{tab:bystage}
\begin{tabular}{llrrrrr}
\toprule
Model & Stage & Action & Crit. ev. & Order & Claim prec. & Overclaim \\
\midrule
\code{gpt-4.1-mini} & Stage 1 & 0.639 & 0.292 & 0.341 & 0.433 & 422 \\
\code{gpt-4.1-mini} & Stage 2 & 0.414 & 0.128 & 0.111 & 0.285 & 506 \\
\code{gpt-4.1-mini} & Stage 3 & 0.345 & 0.164 & 0.151 & 0.382 & 262 \\
\code{gpt-5.4-mini} & Stage 1 & 0.791 & 0.574 & 0.579 & 0.559 & 215 \\
\code{gpt-5.4-mini} & Stage 2 & 0.800 & 0.754 & 0.540 & 0.597 & 213 \\
\code{gpt-5.4-mini} & Stage 3 & 0.802 & 0.779 & 0.540 & 0.594 & 223 \\
\code{gpt-5.5 low raw} & Stage 1 & 0.937 & 0.908 & 0.897 & 0.658 & 374 \\
\code{gpt-5.5 low raw} & Stage 2 & 0.915 & 0.908 & 0.857 & 0.620 & 415 \\
\code{gpt-5.5 low raw} & Stage 3 & 0.903 & 0.903 & 0.849 & 0.643 & 446 \\
\bottomrule
\end{tabular}
\end{table*}

この結果の解釈では注意が必要である．Stage 3で性能が維持されたことは，alert summaryが不要であることを意味しない．Stage 3でもCBC EDR/NGAV telemetryは残っており，モデルはそこからprocess，registry，networkなどの証跡へ到達できる．正しい主張は，alert summary rowsが唯一の復元経路ではなく，十分に強いモデルではnon-alert telemetryからも復元可能な場合がある，ということである．

\subsection{Scenario別結果}

表\ref{tab:scenario}にscenario group別結果を示す．明示的実行チェーンは，少なくとも強いモデルにおいて，証跡回収と順序構成が最も安定した成功領域であり，\code{gpt-5.4-mini}でevidence recall 0.824，\code{gpt-5.5 low raw}で0.975に達した．これらはプロセス名，コマンドライン，ネットワーク接続，レジストリ操作が直接観測されるため，gold behaviorとログ証跡を対応づけやすい．scenario groupは，明示的実行チェーンが\code{network\_service\_behavior}と\code{command\_shell\_execution}，複数段ツールチェーンが\code{collection\_or\_tool\_invocation}と\code{script\_execution\_chain}，意味解釈型が\code{persistence\_registry\_run\_key}に対応する．

\begin{table*}[t]
\centering
\scriptsize
\caption{Scenario group別結果}
\label{tab:scenario}
\begin{tabular}{llrrrrrr}
\toprule
Model & Scenario & Chains & Action & Crit. ev. & Order & Claim prec. & Overclaim \\
\midrule
\code{gpt-4.1-mini} & explicit & 15 & 0.501 & 0.237 & 0.243 & 0.351 & 734 \\
\code{gpt-4.1-mini} & multi-step & 7 & 0.428 & 0.158 & 0.176 & 0.382 & 416 \\
\code{gpt-4.1-mini} & semantic & 1 & 0.506 & 0.148 & 0.167 & 0.444 & 40 \\
\code{gpt-5.4-mini} & explicit & 15 & 0.865 & 0.824 & 0.764 & 0.640 & 345 \\
\code{gpt-5.4-mini} & multi-step & 7 & 0.746 & 0.602 & 0.431 & 0.469 & 288 \\
\code{gpt-5.4-mini} & semantic & 1 & 0.642 & 0.481 & 0.333 & 0.719 & 18 \\
\code{gpt-5.5 low raw} & explicit & 15 & 0.978 & 0.975 & 0.972 & 0.609 & 787 \\
\code{gpt-5.5 low raw} & multi-step & 7 & 0.851 & 0.839 & 0.796 & 0.673 & 421 \\
\code{gpt-5.5 low raw} & semantic & 1 & 0.988 & 0.889 & 0.889 & 0.802 & 27 \\
\bottomrule
\end{tabular}
\end{table*}

複数段ツールチェーンは難しい．\code{gpt-5.4-mini}ではaction recall 0.746に対してorderが0.431，claim precisionが0.469まで下がる．これは，個々の観測行を拾う能力と，複数ツールをまたぐ主系列として順序づける能力が別問題であることを示す．\code{gpt-5.5 low raw}はmulti-stepでもorder 0.796まで改善するが，overclaimは421残る．意味解釈型はDiscord Run Keyの1 chain，9 scored rowsのみであり，広いカテゴリ結論ではなく，レジストリ永続化のcase studyとして扱う．

\subsection{コストと実行上の含意}

表\ref{tab:cost}にコスト効率を示す．\code{gpt-4.1-mini}と\code{gpt-5.4-mini}はlocal ledgerとlocal price estimateに基づく集計であり，\code{gpt-5.5 low raw}はlocal ledgerの\code{call\_total\_usd}がゼロ記録だったため，公開パッケージ内のcost audit noteに基づく保守的推定である．\code{gpt-5.4-mini}は\code{gpt-4.1-mini}より総コストはわずかに高いが，復元精度の改善が大きいため，cost per evidence hitとcost per order hitはむしろ低い．一方，\code{gpt-5.5 low raw}は最も高い精度を示すが，207 runの保守的推定コストは約93.52ドルであり，cost per evidence hitとcost per order hitは大きく増える．

\begin{table}[t]
\centering
\scriptsize
\caption{コスト効率の概算（4.1/5.4はlocal ledger推定，5.5はprogress側再構成推定）}
\label{tab:cost}
\begin{tabular}{lrrrr}
\toprule
Model & Total & Avg/run & /evidence & /order \\
\midrule
\code{gpt-4.1-mini} & \$3.64 & \$0.0176 & \$0.0320 & \$0.0479 \\
\code{gpt-5.4-mini} & \$4.35 & \$0.0210 & \$0.0106 & \$0.0208 \\
\code{gpt-5.5 low raw} & \$93.52 & \$0.4518 & \$0.1764 & \$0.2851 \\
\bottomrule
\end{tabular}
\end{table}

このため，本実験の運用上の結論は，単純に最高精度モデルを選ぶことではない．構造化出力を満たし，低コストで安定した改善を示す\code{gpt-5.4-mini}が現時点の最も現実的な選択肢である．\code{gpt-5.5 low raw}は，raw textからのsalvage採点では強いが，出力契約違反と高コストを考えると，そのまま大規模実験や実運用へ展開するには制約が大きい．なお，\code{gpt-5.5 low raw}の金額はcost audit noteに基づく再構成推定であり，請求額そのものとしては扱わない．

\end{document}
